\title{\bfseries Who Profits When a New Kind of Offering Appears?\\
\large Evidence from 3.1 Million US Trademark Registrations}
\author{Eric Silver\thanks{Independent researcher, Pittsburgh, PA, USA.
ORCID 0000-0003-3351-1109. \texttt{eric.silver@aporal.com}.}}
\date{October 2026}

\begin{document}
\maketitle
\begin{abstract}
\noindent A new kind of offering creates value; whether its creator keeps
that value depends on how easily it is imitated and on who holds the
assets needed to sell it \citep{Teece1986}. We measure both sides on one
outcome for every entrant: whether a registered US trademark's goods are
still in commercial use at the sworn five-year declaration. Each of 3.1
million registrations from 2002--2018 is scored on two positions of its
goods-and-services language within its industry: atypicality, how unusual
the language is, and lead, whether the industry later moved toward it.
Unusual offerings create value that lasts: the most unusual filing in a
class and year is still in use 5.4 points more often than the most
typical, and renewed at year ten 4.3 points more often. Early offerings
lose it: the most leading is in use 2.2 points less often than the most
lagging. Who keeps the value follows the appropriability account. Patents
double the advantage of unusual language and do nothing for the cost of
being early. Growth in entry and the geographic concentration of rivals
halve the advantage. The cost of leading falls on filers without counsel
and on platform businesses, and it is largest in a dozen industries, from
tobacco to technology services, where it lasts to renewal.

\medskip
\noindent\emph{Keywords:} profiting from innovation; first-mover
advantage; appropriability; complementary assets; trademarks; product
survival.
\end{abstract}

\section{Introduction}\label{sec:intro}

A new kind of offering creates value, and the firm that brings it to
market often does not keep it. Where an offering is easy to imitate and
the assets needed to sell it belong to others, imitators and the holders
of those assets capture what the innovator created; where it can be
protected, or the innovator holds the assets itself, the innovator keeps
it \citep{Teece1986}. The question of who profits is the one beneath a
longer argument about entry order.

\emph{Pioneers keep the markets they open.} In surveys of consumer-goods
businesses, market pioneers held larger shares than early followers, and
early followers larger than late entrants \citep{RobinsonFornell1985}. A
pioneer can stay ahead in technology and learning, take scarce inputs
before others arrive, and hold customers who would pay to switch
\citep{LiebermanMontgomery1988}; when a new industry is later winnowed,
the early entrants are the likeliest to survive the shakeout
\citep{KlepperMiller1995, Klepper1997}.

\emph{Most pioneers fail, and followers take what they proved.} The
histories of fifty product categories showed that about half of the true
first entrants failed and that the firms remembered as pioneers had
usually entered later \citep{GolderTellis1993}; in consulting, rivals
copy new offerings routinely \citep{SemadeniAnderson2010}. Whether moving
first pays depends on the resources the entrant brings
\citep{LiebermanMontgomery1998} and on how fast technology and demand
change \citep{SuarezLanzolla2007}.

\emph{Creating a new category pays.} A practitioner literature holds that
firms which create uncontested market space, often by combining
attributes from different industries, grow faster and earn more than
those that compete in existing categories \citep{KimMauborgne2004}.

These claims are argued from cases, most of them firms that survived to
be studied. The trademark record covers every entrant, including those
that failed, and dates what each one offered. A US trademark registration
carries a description of the goods or services the mark covers. Between
the fifth and sixth year after registration the owner must swear that the
mark is still in commercial use on those goods, or the registration is
cancelled; the declaration therefore records whether the product was
still being sold. We score each registration's description on two
positions within its industry \citep{SilverCorpus2026}:
\emph{atypicality}, how unusual the language is against the filings
before and after it, and \emph{lead}, whether the industry's later
filings moved toward it. Atypicality measures a new kind of offering in
its industry: under a theme model fitted across all industries, it
largely records language that brings to an industry what others describe.
Lead measures arriving where the industry was about to go.

Four results follow, organized by the appropriability account. First,
unusual offerings create value that lasts, and early ones lose it: the
most unusual filing in its class and year is still in use 5.4 points more
often than the most typical and the most leading 2.2 points less often
than the most lagging, and both effects hold at the ten-year renewal.
Second, imitation erodes the value: the advantage of unusual language
halves where a class's filing volume grows fast and where the filers of
the same theme concentrate in place. Third, protection keeps it: owners
who patent keep 7.6 points of the advantage against 3.7 for owners who do
not, and patents do nothing for the cost of being early. Fourth, the cost
of being early falls where the entrant lacks complementary assets: on
filers without counsel, on platform businesses, and in a dozen
industries where it lasts to renewal.

% PENDING (agents): who collects within a wave (pioneers, followers,
% incumbents, acquirers); product or firm (owner fixed effects); crowding
% dose-response; industry structure (record-based five-forces and Teece
% scores; blind expert ratings). Add a fifth result paragraph once in.

Section~\ref{sec:data} describes the registrations, the two positions
and the factors. Section~\ref{sec:headline} reports the value created
and lost. Section~\ref{sec:magnitudes} compares every factor on one
scale. Section~\ref{sec:blue} reports what erodes the value of unusual
offerings, and Section~\ref{sec:protect} what lets the creator keep it.
Section~\ref{sec:collect} asks who collects within a wave of entry, and
Section~\ref{sec:structure} whether industry structure predicts where
early entrants lose. Section~\ref{sec:geo} builds the geographic measure
used throughout, Section~\ref{sec:industries} maps the results by
industry, and Section~\ref{sec:discussion} returns to the question of
who profits.
